% concatenated from topicgfib-v41.tex
%\documentclass[12pt]{article}
\documentclass[12pt,reqno]{amsart}
\renewcommand{\baselinestretch}{1.0}
\newenvironment{pf}{\noindent\textbf{Proof.}\quad}{\hfill{$\Box$}}
\usepackage[colorlinks=false]{hyperref}
%\usepackage[pdftex]{hyperref}
\usepackage[dvips]{graphicx}
\usepackage[usenames,dvipsnames]{color}
\usepackage{epsfig}
\usepackage{rotating}
\usepackage{amssymb}
\usepackage{amsmath,amsfonts}
%\usepackage[boxed]{algorithm}
\usepackage[latin1]{inputenc}
\usepackage{verbatim}
\usepackage[tableposition=top,font=small,labelfont=bf,format=hang]{caption}
\usepackage{booktabs}
\newtheorem{defi}{Definition}
\newtheorem{lem}{Lemma}
\newtheorem{thm}{Theorem}
\newtheorem{prop}{Proposition}
\newtheorem{conj}{Conjecture}
\newtheorem{rem}{Remark}
\newtheorem{coro}{Corollary}
\newtheorem{ex}{Example}
%\newtheorem{ex}[thm]{  Example}
\newcommand{\bbF}{\mathbb{F}}
\newcommand{\bbZ}{\mathbb{Z}}
\newcommand{\F}{\mathbb{F}}
\newcommand{\Z}{\mathbb{Z}}
\newcommand{\Q}{\mathbb{Q}}
\newcommand{\R}{\mathbb{R}}
\newcommand{\N}{\overline{N}}
\newcommand{\ep}{\varepsilon}
\newcommand{\eps}[1]{{#1}_{\varepsilon}}%
\newcommand{\vc}{{\bf c}}
\newcommand{\vd}{{\bf d}}
\newcommand\tab[1][1cm]{\hspace*{#1}}
\setcounter{MaxMatrixCols}{20}
\newcommand{\RR}{\mathbb{R}}
%\newcommand{\qed}{\hfill $\Box$ \\}
\newcommand{\rvline}{\hspace*{-\arraycolsep}\vline\hspace*{-\arraycolsep}}
%\newenvironment{proof}{{\em Proof:}}
%\newenvironment{proof}{\noindent\textbf{Proof.}\quad}

%\rule{1ex}{1ex}\\}

\setcounter{page}{1}

\DeclareMathOperator{\wt}{wt}
\renewcommand{\vec}[1]{\boldsymbol{#1}} % Uncomment for BOLD vectors.

%opening
\title{ Second order Recurrences, quadratic number fields and cyclic codes}
\thanks{The work of Minjia Shi was supported in part by the National Natural Science Foundation of China under Grant 12471490.
The work of Jon-Lark Kim was supported in part by the BK21 FOUR (Fostering Outstanding Universities for Research) funded by the Ministry of
Education(MOE, Korea) and National Research Foundation of Korea(NRF)
under Grant No. 4120240415042 and by Basic Science Research Program through the National Research Foundation of Korea(NRF) funded by the Ministry of Science and ICT under Grant No. RS-2025-24534992.}
%\author{ Minjia Shi, Xuan Wang, Zakariae Bouazzaoui, \\ Jon-Lark Kim, Patrick Sol\'e}
\author{Minjia Shi}
\address{Key Laboratory of Intelligent Computing Signal Processing, Ministry of Education,
	School of Mathematical Sciences, Anhui University, Hefei 230601, China}
\email{smjwcl.good@163.com}
\author{Xuan Wang}
\address{Key Laboratory of Intelligent Computing Signal Processing, Ministry of Education,
	School of Mathematical Sciences, Anhui University, Hefei 230601, China}
\email{wang\_xuan\_ah@163.com}
\author{Bouazzaoui Zakariae}
\address{\'Ecole Sup\'erieure de l'Education et de la Formation. Laboratoire de mathématiques et applications FSO-Maroc}
\email{z.bouazzaoui@ump.ac.ma}
\author{Jon-Lark Kim}
\address{Department of Mathematics and Institute for Mathematical and Data Sciences, Sogang University, Seoul, South Korea}
\email{jlkim@sogang.ac.kr}
\author{Patrick Sol\'e}
\address{I2M, Aix Marseille Univ, CNRS,  Marseille, France}
\email{patrick.sole@telecom-paris.fr}

\date{}

\begin{document}
\maketitle
\begin{abstract}
 Wall-Sun-Sun primes (shortly WSS primes) are defined as those primes $p$  such that the period of the Fibonacci recurrence is the same modulo
 $p$ and modulo $p^2.$ This concept has been generalized recently to certain second order recurrences whose characteristic polynomials admit as a zero the principal unit of $\Q(\sqrt{d}),$
 for some integer $d>0.$ Primes of the latter type we call $WSS(d).$ They correspond to the case when $\Q(\sqrt{d})$ is not $p$-rational. For such a prime $p$
 we study the weight distributions of the cyclic codes over $\F_p$ and $\Z_{p^2}$ whose
 check polynomial is the reciprocal of the said characteristic polynomial. Some of these codes are MDS (reducible case) or NMDS (irreducible case).
\end{abstract}

\noindent {\bf Keywords:} WSS primes, quadratic number fields, $p$-rational number fields, second order recurrences, cyclic codes, MDS codes, NMDS codes\\
{\bf MSC (2020):} 11B39, 11B50, 11R11, 94B15
%%%%%%%%%%%%%%%%%%%%%%%%
\section{Introduction}

The famous Fibonacci sequence is defined by the linear recurrence
$$F_{n}=F_{n-1}+F_{n-2},\; n\ge 2,$$
with initial conditions $F_0=0,\,F_1=1.$
This is a periodic sequence modulo every positive integer $m$, its period $\pi(m)$ is called the Pisano period.
Formulas and properties of $\pi(m)$ are well known since Wall's paper in $1960$ \cite{W}. That paper renewed interest in the subject
and motivated studies of Pisano periods, which were later related to other objects and generalized to other classes of second order recurrence
sequences. A natural question was asked by Wall in his paper: Can there be a prime $p$ such that $\pi(p)=\pi(p^2)?$
It is known that up to $10^{14}$, there are no such primes (cf. \cite{MR}). Still, using heuristics and probabilistic arguments, some authors conjecture the existence of infinitely many primes $p$ satisfying $\pi(p)=\pi(p^2)$ \cite{EJ,Kl}. It was shown by the Sun brothers in \cite{SS} that the existence of these primes is related to Fermat's Last Theorem, the reason why we call them Wall-Sun-Sun primes, or WSS primes for short.\\
The question of Wall can naturally be asked for other types of second order recurrence sequences. In \cite{B,B2}, this was considered for generalized Fibonacci sequences of the form $$F_{n}=-AF_{n-1}-BF_{n-2},\; n\ge 2,$$
whose characteristic polynomial is $X^2+AX+B$, with $A,B \in \Z$ being associated with a given real quadratic field. The question of Wall for these sequences is related to certain deep arithmetic properties of real quadratic fields. More precisely, let $d>0$ be a fundamental discriminant and let $\varepsilon_d$ be the fundamental unit of $\Q(\sqrt{d})$. Then $\varepsilon_d =\frac{a+b\sqrt{d}}{2}$ when $(a,b)$ represent the smallest solution of the Pell equation $x^2-dy^2=\pm4$. The sequence in \cite{B,B2} is then given by $$F_{n+2}=aF_{n+1}-\frac{a^2-db^2}{4}F_n,$$ where $X^2-aX+\frac{a^2-db^2}{4}$ is the minimal polynomial of $\varepsilon_d$. Write $k(m)$ for the period of this sequence modulo $m>1$ and say that a prime $p$ is a $WSS(d)$ prime if we have $k(p^2)=k(p)$. Then for a large enough prime $p$, we have $k(p^2)=k(p)$ exactly when the quadratic field $\Q(\sqrt{d})$ is $p$-rational \cite{B,B2} (see Section \ref{p-rational fields} below for details about $p$-rational fields). In this setting, for the classical Fibonacci sequence, that is $d=5$, there are no known $WSS(5)$ primes. However, for other values $d>5$, one can find many $WSS(d)$ primes \cite[Table]{B}.
\vspace{1em}

To allow more generality, we shall say that a prime is $WSS(d)^*$ if the characteristic polynomial in a general form discussed above has discriminant $d$ up to square factors, but not necessarily $\varepsilon_d$ as a zero.
With every second order recurrence, can be attached a cyclic code of dimension two, whose codewords are the period window of the sequence when the initial conditions vary.
The weight distribution of cyclic codes of dimension two is studied over finite fields in \cite{SZS,V1,V2} and over rings in \cite{SHS,SLS}. The methods in \cite{SHS,SLS,SZS}
are based on a local Binet-like formula; while the technique in \cite{V1,V2} is a more classical coding theoretic approach.
Their salient property is that they are one-weight code or
two-weight codes, a fact which ties them up with strongly regular graphs in the projective case \cite{BM}. Some of our codes are MDS (reducible case) or  near MDS (shortly NMDS) (irreducible case). NMDS codes is an important class of codes introduced in \cite{DL} that resembles MDS codes.
For some primes $p,$ we construct in Proposition 5,  NMDS $[2p+2,2,2p]$ codes, that are optimal for the length.

The material is organized as follows. The next section contains prerequisite notions. Section 3 develops the connection between WSS primes and quadratic number fields. Section 4 shows that, given a prime $p>0$, there is an infinitude of integers $d >0$ such that $p$ is $WSS(d)^*.$ Section $5$ considers the cyclic codes over $\F_p$ and $\Z_{p^2}$ of length the common period when $p$ is $WSS(d).$ Section $6$ concludes the article.
%%%%%%%%%%%%%%%%%%%%%%%%%%%%%%%%%%%%%%%%%
\section{Preliminary}
\subsection{Codes}
We assume that the reader has some basic familiarity with Coding Theory as can be gained from the first chapter of \cite{MacSlo} or \cite{HuffmanPless}.

The {\bf Hamming weight} of a vector $\vec{u}$ over a finite field $F$ and of length $n$ is the number of all nonzero coordinates of $\vec{u}$, denoted by $\wt(\vec{u})$.
A {\bf code} of length $n$ over a finite field $F$ is a vector subspace of $F^n.$
Its elements are called {\bf codewords}.
Its {\bf dimension} is denoted by $\kappa,$ and its minimum weight or {\bf minimum distance} by $d.$
The three parameters are written compactly as $[n,\kappa,d]$ or $[n,\kappa,d]_q$ if $F=\mathbb F_q$, the finite field of order $q,$ where $q=p^r$ ($p$ a prime and $r\ge1$).
An $[n,\kappa,n-\kappa+1]$ is called {\bf MDS}.
An $[n,\kappa,n-\kappa]$ is called almost-MDS, simply {\bf AMDS}
and near-MDS, simply {\bf NMDS} if its dual is also AMDS.
A {\bf generator matrix} $G$ of the code is any $\kappa$ by $n$ matrix whose rows span the code.

A code is {\bf projective} if any pair of columns of $G$ is linearly independent.
A code is {\bf two weight} if it has only two nonzero weights $w_1,\,w_2$ with $w_1 < w_2$.  This is denoted
by $[n,k,\{w_1,\,w_2\}].$ It was shown by Delsarte that if the code is projective then $w_2-w_1$ is a power of the
charcteristic $p$ of $F.$
If  $m>1$ is the largest factor of $w_2-w_1$ that is coprime to $p$,
 then by \cite{K} the code is an $m$-fold repetition of a projective code. (See \cite{K} for a precise definition of the $m$-fold repetition).

Considering codes over a finite commutative ring $R$ we call linear codes of length $n$ over $R$ the
submodules of $R^n.$ We use the notation $(n,M,d)$ for codes with $M$ the size and $d$ the minimum distance of the code.

In the present paper we will use $F=\F_p$ and $R=\Z_{p^2}$ for some prime $p.$
The number of codewords of weight $w$ is denoted by $A_w$ both in the field and ring case. The sequence $(A_w)_{w=0}^n$ is called the {\bf weight distribution} of the code.

\subsection{Cyclic codes}

A {\bf cyclic code} of length $n$
over a finite field $F$ is an ideal of the polynomial ring $F[x]/(x^n-1).$ It can be shown that this ideal is of the form $(g),$ for some monic divisor $g$ of $x^n-1.$ The polynomial $g$ is called the
{\bf generator polynomial}  and the  {\bf check polynomial} is defined as $\frac{x^n-1}{g}.$
Similarly
a cyclic code of length $n$
over a finite ring $R$ is an ideal of the polynomial ring $R[x]/(x^n-1).$ For a free code we define similarly its generator and check polynomials.
Given a second order recurrence $$F_{n}=-AF_{n-1}-BF_{n-2},\; n\ge 2,$$ of period $T$
over $R$ there is a $[T,2]$ cyclic code $C$ that can be constructed from it with codewords
$$ (F_0,\dots,F_{T-1})$$
and initial conditions arbitrary. Thus
$$C=\{(F_0,\dots,F_{T-1}) \mid F_0,F_1 \in R\}.$$
It can be shown that the reciprocal polynomial of the characteristic polynomial of the recurrence is, up to normalization, the check polynomial of $C.$

\subsection{Basic Number theory} \label{sub-number-theory}
The {\bf Legendre symbol} $(\frac{.}{p})$ is defined for an odd prime $p$ by the relation
$$\left(\frac{x}{p}\right)=
                \begin{cases}
      1, & \text{if}\ x \in Q \\
      0 ,& \text{if}\ x=0\\
      -1, & \text{otherwise}
    \end{cases}
 $$
 where $Q$ is the set of quadratic residues modulo $p$ defined by
 $$Q=\{x^2 \mid x \in \F_p^\times\}.$$

 In particular a quadratic $t^2+ut+v$ is reducible over $\F_p,$ for some odd prime $p$  if and only if its {\bf discriminant}
 $\Delta_p=u^2-4v^2 \in \F_p$ is such that $(\frac{\Delta_p}{p})\in \{0,1\}.$\\ \ \

 A {\bf number field} is a finite extension of the rationals. It is {\bf quadratic} if the degree of the extension is two. It is {\bf real quadratic} if it is of the form $\mathbb{Q}(\sqrt{d})$ for some square-free integer $d>1.$ The {\bf class number} of a number field $F$ is the order of its class group; that is the group of fractional ideals of $F$ modulo principal ideals. A field extension $F/K$ is said to be Galois if it is normal (every irreducible polynomial over $K$ that has one root in $F$ splits completely in $F$) and separable (all extensions of ${\Q}$ are separable). The Galois group of $F/L$ is the group of automorphisms of $F$ that fixes $K$ pointwise.
%%%%%%%%%%%%%%%%%%%%%%%%%%%%%%%%%%%%%%%%%%%%%%%%%%%%%%%%%%%%%%%%횄쨔

%\subsection{Number fields}
%%%%%%%%%%%%%%%%%%%%%%%%%%%%%%%%%%%%%%%%%%%%%%%%%%%%%%%%%%%%%%%%횄쨔
\section{Real quadratic $p$-rational fields}\label{p-rational fields}



In this section, we briefly introduced $p$-rational real quadratic fields, where $p$ is an odd prime number, and their precise connection to periods of Fibonacci numbers \cite{B,B2}.

Let $F$ be a number field and let $p$ be a rational prime. A $p$-adic prime $\mathfrak{p}$ of $F$ is an ideal of the ring of integers of $F$ above the prime $p$. Let $S_p$ be the set of $p$-adic primes of $F$. A $p$-extension $K$ of $F$ is an extension of $F$ of degree a $p$-power, and it is called unramified outside $S_p$ if the only possible prime ideals that are ramified in $K/F$ are those in $S_p$. The compositum $F_{S_p}$ of all such $p$-extensions is called the maximal pro-$p$-extension of $F$ that is unramified outside $S_p$. Its Galois group $\Gamma_F$ is a pro-$p$-group; a projective limit of finite $p$-groups \cite[Section 1.4]{Se}.

Let $\Phi(\Gamma_F)$ denote the Frattini subgroup
of $\Gamma_F$, which is defined to be the intersection of all closed subgroups of index $p$.
The Frattini quotient of $\Gamma_F$ is the quotient group $\tilde{\Gamma}_{F}:=\Gamma_F/\Phi(\Gamma_F)$.
This is an abelian
group of exponent $p$, which we regard as a vector space over the field ${\F}_p$. A number field $F$ is called $p$-rational if
the ${\F}_p$-dimension of $\Tilde{\Gamma}_{F}$ is exactly $r_2(F)+1$, where $r_2(F)$ is the number of complex embeddings of $F$, see \cite[Theorem and Definition 2.1]{Movahhedi-Nguyen} for several equivalent formulations of the $p$-rationality.


%Following \cite[Theorem and Definition 2.1]{Movahhedi-Nguyen}, The field $F$ is called $p$-rational if the Galois group $G_{S_p}(F)$
%is a free pro-$p$-group
%A general number field $F$ is called $p$-rational if the Galois group $G_p(F)$ of the maximal pro-$p$-extension of $F$ that is unramified outside the set of $p$-adic primes is a free pro-$p$-group (c.f. \cite[Theorem and Definition 2.1]{Movahhedi-Nguyen} with other several equivalent formulations of the $p$-rationality), for details about pro-$p$-groups see \cite{Se}.

The structure of the Galois group $\Gamma_F$ reflects arithmetic properties of the base field $F$ at the prime $p$.
More precisely, let $\Gamma_{F}^{ab}$ be the maximal abelian quotient of $\Gamma_F$; that is the Galois group of the maximal
abelian extension contained in $F_{S_p}$. We have the following isomorphism: $$\Gamma_{F}^{ab}\cong\mathbb{Z}_{p}^{1+r_2(F)+\delta_p}\times T_F,$$
where $\delta_p$ is Leopoldt's defect (\cite[Corollary 10.3.7]{NSW}), $T_F$ is a finite $p$-group representing the $\mathbb{Z}_p$-torsion and ${\Z_p}$ is the ring of $p$-adic integers.
It follows that $F$ is $p$-rational
if and only if $\delta_p=0$ and $T_F$ is trivial, where the vanishing of $\delta_p$ means that Leopoldt's conjecture is true \cite[Section 10.3]{NSW}.
When the number field $F$ is totally real, Leopoldt's conjecture for $F$ at $p$ is expressed by the non vanishing
of the $p$-adic regulator $R_p(F)$ (c.f. \cite[10.3.4 and 10.3.5]{NSW}). Furthermore,
a formula linking the order of $T_F$ with $R_p(F)$ was
proved by Coates in \cite[Lemma 8]{Co}. When $F$ is a real quadratic field, this formula reads
$$|T_F|\sim_p h_F R_p(F)/p^{1/e},$$
where $\sim_p$ means that both sides have the same $p$-adic valuation, $h_F$ is the class number of $F$ and $e := 2$ or $1$ according as $p$ ramifies in $F$ or not.
If $\varepsilon$ is the fundamental unit of $F$, then we have
$R_p(F)=\log_p(\varepsilon)$, where $\log_p$ is the $p$-adic logarithm. We use the fact that the $p$-adic valuation of $R_p(F)$ is the same as that of $(\varepsilon^{Nv-1}- 1)$, where $v$ is a $p$-adic prime of $F$ and $Nv$ is its absolute norm, to relate the $p$-rationality of $F$ in an elementary way to the divisibility by $p^2$ of certain associated Fibonacci numbers. The starting point is the following characterization:

%In particular, if $F$ is $p$-rational then it satisfies Leopoldt's conjecture \cite[Section 10.3]{NSW}.


\begin{prop}(\cite[Proposition 4.1]{Greenberg})
	Let $F$ be a real quadratic number field. Suppose either $p\geq 5$ or $p =3$, and is unramified in
	$F$. Then $F$ is $p$-rational precisely when the following two conditions are satisfied:
	\begin{enumerate}
		\item[(i)] The prime $p$ does not divide the class number $h_F$,
		\item[(ii)] The fundamental unit of $F$ is not a $p$-th power in the completion $F_v$ of $F$ at a p-adic place $v$.
	\end{enumerate}
\end{prop}

The class number $h_F$ being finite, the most important property to study in the real quadratic case is the behavior of the fundamental unit
locally at $p$. Computations suggest that primes $p$ for which a number field is not $p$-rational are rare in general, for example $\mathbb{Q}(\sqrt{5})$
is $p$-rational for all primes up to $10^{14}$. Greenberg proved an elementary criterion to check the $p$-rationality of the field $\mathbb{Q}(\sqrt{5})$
by means of congruences of Fibonacci numbers \cite[Corollary 4.5]{Greenberg}. This inspired a general phenomenon relating $p$-rationality to periods of Fibonacci numbers.

Let $d>0$ be a fundamental discriminant. Denote by $\varepsilon_d$
and $h_d$ the fundamental unit and the class number of the field
${\Q}(\sqrt{d})$ respectively, and let $\overline{\varepsilon}_d$ be the conjugate of $\varepsilon_d$. We associate to the field ${\Q}(\sqrt{d})$ a Fibonacci
sequence
$F^{(\varepsilon_d+\overline{\varepsilon}_d,-\varepsilon_d \overline{\varepsilon}_d)}=(F_n)_{n\geq0}$
defined by $F_0=0$, $F_1=1$ and the recursion formula
\begin{equation}\label{recurrence for d}
F_{n+2}=(\varepsilon_d+\overline{\varepsilon}_d)F_{n+1}-\varepsilon_d \overline{\varepsilon}_d F_{n},\; \text{for\;$n\geq0$}.
\end{equation}

Let $p\geq5$ be a prime number such that $p\nmid
(\varepsilon_{d}-\overline{\varepsilon}_{d})^2 h_d$. Then,

\begin{equation*}\label{p-rational and congruences}
{\Q}(\sqrt{d})\; \hbox{is}\; p\hbox{-rational if and only
	if}\; F_{p-(\frac{d}{p})}\not\equiv0\pmod{p^2},
\end{equation*}
where $(\frac{d}{p})$ denotes the Legendre symbol (see \cite[Theorem 3.4]{B2}).
The proof of this result uses the equivalence

\begin{equation*}\label{regulator}
{\Q}(\sqrt{d})\;\hbox{is}\;
p\hbox{-rational}\;\Leftrightarrow\; \log_{p}{(\varepsilon_d)}\not\equiv0\pmod{p^{2}},
\end{equation*}
where $\varepsilon_d$ is a fundamental unit of $K$ and $\log_p$ is
the $p$-adic logarithm.

\begin{thm}\label{periods and p-rationality}(\cite[Theorem 1.1]{B})
	Let $d>0$ be a fundamental discriminant. For every odd prime number
	$p$ such that $p\nmid
	(\varepsilon_d-\overline{\varepsilon_d})^{2}h_d$, the quadratic field ${\Q}(\sqrt{d})$ is
	$p$-rational precisely when $k(p)\neq k(p^{2})$.
\end{thm}

Let us consider the polynomials of the form $P_{d}(X)=X^2-aX+\frac{a^2-db^2}{4}$. The splitting field
of the polynomial $P_d$ is the quadratic field $\mathbb{Q}(\sqrt{d})$. The unit defined by $P_d$ is then a power of the fundamental
unit of $\mathbb{Q}(\sqrt{d})$. By Theorem 3.4 of \cite{B2}, if $k(p)= k(p^2)$ then $\mathbb{Q}(\sqrt{d})$ is necessarily not $p$-rational, under the assumption that $p\nmid(\varepsilon_d-\overline{\varepsilon_d})^{2}h_d$.
\vspace{1em}

We consider the quadratic field $\mathbb{Q}(\sqrt{5})$, for which no $WSS(5)$ primes are known to exist.
By Corollary $4.3$ of \cite{B2} we have the following elementary characterization of the $p$-rationality of
$\mathbb{Q}(\sqrt{5})$ for primes $p\neq5$. For $0<i<p$, we denote by $\frac{1}{i}$ the inverse
of $i$ modulo $p$ in $\mathbb{F}_{p}^*$. For a fundamental discriminant $d$, we set
\begin{equation*}
\alpha_{p}^{d}(i)=\sum_{\begin{array}{c}
	1\leq k\leq \frac{p-1}{2}\\
	k\equiv i\pmod{d}
	\end{array}
}\frac{1}{k}\pmod{p}.
\end{equation*}
For every prime number $p\equiv1\pmod{5}$, we have the equivalence:
\begin{equation*}
\text{$\mathbb{Q}(\sqrt{5})$ is not $p$-rational $\iff$ $\alpha_p^5(1)+2\alpha_p^5(5)\equiv\alpha_p^5(4)-\alpha_p^5(2)\pmod{p}$.}
\end{equation*}









%\section{Questions}
%Let $p$ be a WSS(d) prime.
%Consider the cyclic code $C_1$ (resp. $C_2$) of dimension $2$ with check polynomial $X^2-aX+%{a^2-db^2}{4}$ over $\F_p$ (resp. $\Z_{p^2}$) of length $k(p)=k(p^2).$
%These codes can be either irreducible or reducible depending whether $d$ is a quadratic residue modulo $p$ or not.
%Can the properties of SRGs \cite{BM} be used to derive non existence results for $WSS(d)$ primes ? Does the fact that $C_1$ is both the residue and torsion code of $C_2$ give some constraints on the weight distribution of both codes?
%Constructing $C_1$ and $C_2$ explicitly for the primes $p$ in the table of \cite{B}, and their corresponding SRGs can give some insights into these questions.
%%%%%%%%%%%%%%%%%%%%%%%%%%%%%%%%%%%%%%%%%%%%%%%%
\section{ Inverse problem} \label{sec:inv-prob}
Given a prime $p>5$ we can construct infinitely many $d$'s such that $p$ is WSS$(d)^*.$ This result is proved  in Theorem 1.1 of \cite{QS} by a different method. We need some basic results on polynomials over finite fields. Let $q$ be a prime power. The {\bf order} $e$ of a polynomial $f(x)$ of
$\F_q[x]$ is the smallest integer $e$ such that
$f(x) | x^e-1.$ (This integer is known to exist provided $f(0)\neq 0$).

{\lem \label{basic} The order of $f\in \F_q[x]$ equals the largest order of its roots in the multiplicative group of the splitting field of $f$ over $\F_q.$}

See \cite[Th. 8.27]{LN} for a proof.

We take for granted that the period of a linear recurrence over a finite field equals the period of its characteristic polynomial \cite[Th. 8.28, p.408]{LN}.
It is known that either $k(p^2)=k(p)$ or $k(p^2)=p k(p)$ \cite[Prop. 1, p. 375]{R}. Hence if we show that
$k(p^2)\le k(p)$ we can conclude that $k(p^2)=k(p).$

%%%%%%%%%%%%%%%%%%%%%%%%%%%%%%%%%%%%%%%%%%%%%%%%
\subsection{Reducible cyclic codes}
\subsection{Two roots} \label{sec-order-p-1}
By Fermat, the polynomial $x^{p-1}-1$ splits into linear factors over $\F_p.$ Let $h(x)$ be the product of any two of these factors. Assume that the order of $h(x)$ is $p-1,$
or, by Lemma \ref{basic}, that one at least of its roots has order $p-1.$
The factorization $x^{p-1}-1=h(x)g(x)$ can be lifted into a factorization $x^{p-1}-1=H(x)G(x)$ over $\Z_{p^2}[x]$ with $H,G$ being congruent to $h,g$ respectively
$\pmod{p}.$ This is an immediate application of Hensel lemma \cite{McG}. Note that $H$ is also of order $p-1.$
Any lower order would lead to a contradiction with the order of $h$ by reduction modulo $p.$
Write $H(x)=x^2+ax+b.$ Consider arbitrary integers $A,B$ congruent to $a,b$ respectively $ \pmod{p^2}.$ We compute the discriminant $\Delta=A^2-4B$ of
$X^2+AX+B.$
We can choose $A,B$ such that
$\Delta>0.$
Let $d$ be the square free part of $\Delta.$ Then we claim that $p$ is $WSS(d)$ with $k(p)=k(p^2)=p-1.$

\begin{ex}{\rm
%{\bf Example:}
Let $p=31$ for $d=2.$ When $n=30=k(31)=k(961),$ the cyclic code $C_1$ with check polynomial $h(x)$ is a $[30, 2, 29]_{31}$ MDS code.
}
\end{ex}
%%%%%%%%%%%%%%%%%%%%%%%%%%%%%%%%%%%
\subsection{Double root}
Over $\F_p$ we have $x^p-1=(x-1)^p$ implying that
$(x-1)^2=x^2-2x+1$ divides $x^p-1.$ Since the period of
$F_{n+2}=2F_n-F_{n+1}$ divides $p$, it has to be $p.$
Now the polynomial $x^2-2x+1$ can be lifted by the Hensel lemma
to some quadratic $x^2+ax+b$ that divides $x^p-1$ over $\Z_{p^2}$ (We conjecture  based on numerical evidence that $a=-2,\,b=1.$ This fact is not needed in what follows). We can choose in infinitely many ways integers $A,B$ such that
\begin{itemize}
 \item $x^2+Ax+B \equiv x^2+ax+b \pmod{p^2}$
 \item $\Delta:=A^2-4B>0.$
\end{itemize}
Let $d$ be the square free part of $\Delta.$
Then $p$ is $WSS(d),$ since $k(p)=k(p^2)=p.$\\ \ \

\begin{ex}{\rm
 Consider the quadratic $x^2+10x+1 \in \Z[x],$ of discriminant $96=6\times 16.$ Over $\F_3$ and $\Z_9$ it reduces to $x^2+x+1,$ which divides $x^3-1.$ It can be seen that the initial
conditions $F_0=0,\,F_1=1$ of the recurrence
$$F_{n+2}=-(F_n+F_{n+1})$$
yield the sequence of period 3
$$ 0, 1, -1, 0, 1, -1, \dots$$
both over $\F_3$ and $\Z_9.$
The discriminant of $x^2+10x+1 \equiv x^2+x+1 \pmod{9}$ is
$100-4=6\times 16.$ Hence the prime $3$ is $WSS(6)^*.$
}
\end{ex}

\begin{rem}\label{triv}{\rm
 The existence of $\alpha \in \F_p^\times$ of order $p-1$
 entails the existence of $\beta $ of order $D$ for any divisor $D$ of $p-1.$ It is enough to take $\beta=\alpha^{D'}$ with $DD'=p-1.$ This yields, by the same argument as above, the existence
 of a $d$ such that $k(p)=k(p^2)=D.$
 }
\end{rem}

%%%%%%%%%%%%%%%%%%%%%%%%%%%%%%%%%%%%%%%%%%횄쨔
\subsection{Irreducible cyclic codes}\label{subsec-irr-cyc}
%We construct cyclic codes $C_1$ and $C_2$ of length $2p+2.$ over $\F_p$ and $\Z_{p^2},$ respectively.
We need an irreducible quadratic polynomial $h(x)$ over $\F_p$ of order ${2p+2}$.
\begin{prop}
 There exists $\alpha \in \F_{p^2}\setminus \F_p$ satisfying $\alpha^{p+1}=-1.$
 \end{prop}

\begin{proof}
 Such an $\alpha$ exists because $x^{p+1}+1$
divides $x^{2p+2}-1$ which divides itself
$x^{p^2-1}-1$.\\ We can assume that this $\alpha \notin \F_p,$ because the only roots common to that polynomial and to $x^{p-1}-1$ are roots of $x^4-1.$
 This can be seen from $GCD(x^{2p+2}-1,x^{p-1}-1)|x^4-1$ which results in turn from $GCD(2p+2,p-1)| 4$ . Therefore the only roots of
 $x^{2p+2}-1$ in $\F_p$ are roots of $x^4-1.$ Since $2p+2>4$  there are roots of $x^{2p+2}-1$ not in $\F_p.$
\end{proof}

This implies that the order of $\alpha$ in $\F_{p^2}^\times$ is $2p+2.$ We can now construct the $h$ we need.

\begin{prop} \label{prop-order-2p+2}
 Let $\alpha$ be a root as in Proposition 1.
 Let $h(x)=(x-\alpha)(x-\alpha^p).$ Then $h(x)$
 is an ireducible polynomial of
 $ \F_p[x]$ of order $2p+2.$
\end{prop}


 \begin{proof}
  By Frobenius action it is in $\F_p[x]$ and irreducible since its roots are not in $\F_p.$
 It divides $x^{2p+2}-1$ since both $x-\alpha$
 and $x-\alpha^p$ do, and are pairwise coprime.
 The order of $h(x)$ is $2p+2$ since it is the common order of both of its roots.
 \end{proof}


 %If $p=3$ we can take $h(x)=x^2+1$ which divides $x^8-1.$ It can be checked by direct computation that $h(x)$ has order $8.$

 The rest of the argument is similar to the reducible case. Then we claim that $p$ is $WSS(d)^*$ with $k(p)=k(p^2)=2p+2.$

\begin{ex}{\rm
 $p=13$ for $d=2.$ When $n=28=k(13)=k(169),$ we have a non-projective $[28, 2, 24]_{13}$ cyclic code with check polynomial $h(x)$, which is the 4-fold repetition of a $[7,2,6]_{13}$ cyclic code, an MDS code. This
 is consistent with the main result of \cite{K} which claims that two-weight codes whose weights difference is not a power of the characteristic of the field are repetitive.
}
\end{ex}

\begin{rem}{\rm
 The existence of $\alpha \in \F_p^\times$ of order $2p+2$
 entails the existence of $\beta $ of order $D$ for any divisor $D$ of $2p+2$ by the same argument as in Remark 1. This yields, by the same technique as for $D=2p+2,$ the existence
 of a $d$ such that $k(p)=k(p^2)=D.$
 }
\end{rem}
 %%%%%%%%%%%%%%%%%%%%%%%%%%%%%%%%%%%%%%%%%%%%%%%%

\section{Comparing weight distributions}

Let $h(x) \mid x^n-1$ be the check polynomial of the free cyclic code $C_2$ of length $n$ and dimension $2$ over $\mathbb{Z}_{p^2}$, where $\gcd(n,p)=1$.
Let $C_1$ be the cyclic code of length $n$ and dimension $2$ over $\mathbb{Z}_p$ with the check polynomial $\bar{h} \equiv h \pmod{p}$, i.e., $h$ is the Hensel lift of $\bar{h}$.

It is known that the weight enumerator of an MDS code over $\mathbb{F}_q$ is unique (see, e.g., \cite[12, Ch.11, \S 3, Theorem 6]{MacSlo}), which is given by $A_0 = 1$, $A_w = 0$ for $0 < w < d$, and
\[
A_w = \binom{n}{w} (q-1) \sum_{j=0}^{w-d} (-1)^j \binom{w-1}{j} q^{w-d-j}.
\]
We require some facts on the algebraic structure of cyclic codes over $\mathbb{Z}_{p^2}.$
\begin{prop}\cite[Proposition 2]{WS-DM}
	Let $C$ be a cyclic code of length $n$ over $\mathbb{Z}_{p^2}$ with $\gcd(n,p) = 1$.
	Then there exist polynomials $f_1, f_2, f_3, \in \mathbb{Z}_{p^2}[x]$ such that $f_1 f_2 f_3 = x^n-1$ and $C = \langle f_1 f_2, pf_1 \rangle$.
	Moreover, we have $pC = \langle pf_1 \rangle$.
	In particular, if $C$ is free, then $f_2 = 1$.
\end{prop}
 We are now ready for the main result of this section.
\begin{thm} \label{thm-wd}
	Assume that $C_1$ is a two-weight MDS code whose length $n$ is at most $p+1$.
	Then
	\[ A_{n-1}(C_1) = n(p-1), \quad A_n(C_1) = (p-1)(p-n+1), \]
	and
	\[ A_{n-1}(C_2) = n(p^2-1), \quad A_n(C_2) = (p^2 - 1)(p^2 - n+1). \]
\end{thm}
\begin{proof}
	Since $C_1$ is a two-weight MDS code, then $w_1 = n-1$, $w_2 = n$, and
	\[ A_{n-1}(C_1) = n(p-1), \quad A_n(C_1) = (p-1)(p-n+1). \]
	Consider the map $\pi: C_2 \rightarrow C_1$ defined as $\pi(\vec{c}) = \vec{c} \pmod{p}$ for each $\vec{c} \in C_2$.
	So $\pi$ is homomorphism.
	It is clear that for each $\vec{c} \in C_2$, it can be uniquely expressed by $\vec{c}_1 \oplus p \vec{c}_2$, where $\vec{c}_1 = \pi(\vec{c}),\vec{c}_2 \in C_1$, and $\oplus$ means the addition operation over $\mathbb{Z}_{p^2}$.
	Obviously, the Hamming weight of $\vec{c}$ is at least $\pi(\vec{c})$, i.e., $\wt(\vec{c}) \geqslant \wt(\pi(\vec{c}))$.
	Thus, we have the following three cases:
	\begin{enumerate}
		\item[(1)] If $\pi(\vec{c}) = \vec{0}$, then $\vec{c} \in pC_2 = \langle pf_1 \rangle$ since $C_2$ is free.
		Hence, there are $n(p-1)$ such $\vec{c}$ of Hamming weight $n-1$, and $(p-1)(p-n+1)$ such $\vec{c}$ of Hamming weight $n$.
		
		\item[(2)] If $\pi(\vec{c}) \neq \vec{0}$ has Hamming weight $n$, then $\vec{c}$ has Hamming weight $n$.
		And there are $(p-1)(p-n+1)$ such $\vec{c}$ of Hamming weight $n$.
		
		\item[(3)] If $\pi(\vec{c}) \neq \vec{0}$ has Hamming weight $n-1$, then $\vec{c}$ has Hamming weight at least $n-1$.
		Note that the row of the generator matrix (in standard form) of $C_2$ has Hamming weight $n-1$.
		We claim that $\vec{c}$ has Hamming weight $n-1$ if and only if $\vec{c}_2 = \lambda \vec{c}_1 \pmod{p}$ for some $\lambda \in \mathbb{Z}_p$.
		For otherwise, we have the following two cases.
		\begin{enumerate}
			\item If $\vec{c}_2$ has Hamming weight $n$, then it is easy to check that the Hamming weight of $\vec{c}_1 \oplus p \vec{c}_2$ is exactly $n$.
			
			\item If $\vec{c}_2$ has Hamming weight $n-1$ but $\vec{c}_2 \neq \vec{c}_1$ for each nonzero $\lambda \in \mathbb{Z}_p$, then the Hamming weight of $\vec{c}_1 \oplus p \vec{c}_2$ is still exactly $n$.
		\end{enumerate}
		Hence,
		\begin{align*}
			|\{ \vec{c} \in C_2: \wt(\vec{c}) = n-1 \}| & = |\{ (\vec{c}_1, \lambda \vec{c}_1): \vec{c}_1 \in C_1, \wt(\vec{c}_1) = n-1 \}| \\
			& = np(p-1).
		\end{align*}
		the number of codewords $\vec{c}$ of Hamming weight $n-1$ is $n(p-1)$.
	\end{enumerate}
	Therefore, $C_2$ is a two-Hamming-weight code, and
	\[ A_{n-1}(C_2) = n(p-1) + np(p-1) = n(p+1)(p-1) = n(p^2 - 1), \]
	which implies that
	\[ A_{n}(C_2) = p^4 - 1 - A_{n-1} = (p^2 - 1)(p^2 + 1 - n). \]
	We complete the proof.
\end{proof}

\begin{ex}{\rm
    By Magma \cite{magma}, we have found the polynomials such that $C_1$ and $C_2$ have only two Hamming weights $w_1, w_2$ $(w_1 < w_2)$.
    The parameters of such codes are listed in Table \ref{tab-two-weight}, where $a_1 = A_{w_1}(C_1)$, $a_2 = A_{w_2}(C_1)$, and $A_1 = A_{w_1}(C_2)$, $A_2 = A_{w_2}(C_2)$.
    Moreover, ``IRR'' means that the check polynomial is irreducible, and ``RED'' means that the check polynomial is reducible.

    \begin{table}[htbp]
	\caption{Cyclic $[n, 2, \{w_1,w_2\}]_p$ codes  for some WSS($d$) primes $p$ of ~\cite[Table(p. 109)]B} \label{tab-two-weight}
	\begin{tabular}{lllll}
		\toprule
		$(n,d,p)$ & $\{w_1,w_2\}$ & $(a_1,a_2)$ & $(A_1,A_2)$ & poly \\
		\toprule
		$(28,8,13)$ & \{24,28\} & (84,84) & (1176,27384) & IRR  \\
		$(30,8,31)$ & \{29,30\} & (900,60) & (28800,894720) & RED  \\
		$(104,12,103)$ & \{102,104\} & (5304,5304) & (551616,111999264) & IRR   \\
		$(484,13,241)$ & \{480,484\} & (29040,29040) & $(7027680,\approx 3\times 10^9)$ & IRR   \\
		$(8,24,7)$ & \{6,8\} & (24,24) & (192,2208) & IRR \\
		$(174,24,523)$ & \{172,174\} & (45414,228114) & $(23796936,\approx 7\times 10^{10})$ & RED  \\
		\bottomrule
	\end{tabular}
\end{table}
}
\end{ex}


\begin{coro}
	Keep the above notations.
	Then we have $A_{w_i}(C_1)$ divides $A_{w_i}(C_2)$ for $i=1,2$, when $n = p-1$ and $n = (p+1)/2$.
\end{coro}
\begin{proof}
	For each $n$, we always have $A_{n-1}(C_1)$ divides $A_{n-1}(C_2)$.
	As for $A_n$, it is clear that $A_{n}(C_1)$ divides $A_{n}(C_2)$ if and only if
	\[ (p-n+1) \mid (p+1)(p^2 - n + 1) \]
	which is true when $n = p-1$ and $n = (p+1)/2$.
\end{proof}

Let $C$ be a two-weight over $\mathbb{Z}_p$ with nonzero weights $w_1 < w_2$.
According to \cite{K}, if $w_2 - w_1$ is not a power of $p$, then $C$ is non-projective and $C$ is the $\ell$-fold repetition of a smaller two-weight code, where $\ell$ is the largest factor of $w_2 - w_1$ that is coprime to $p$.

\begin{coro}
	Keep the above notations.
	If $n > \ell^2$ and $C_1$ is an $\ell$-fold repetition of a smaller two-Hamming-weight code $C'_1$ which is MDS, then $C_1$ and $C_2$ have the same Hamming weights.
\end{coro}
\begin{proof}
	Let $n = m \ell$.
	Then $C_2$ has nonzero Hamming weights $\ell(m-1) = n - \ell$ and $m \ell = n$.
	Since $C_1$ is an $\ell$-fold repetition, then by a suitable column permutation, $C_1^{(i)}$ can be MDS, where
	\[ C_1^{(i)} = (c_{(i-1)m + 1}, \dots, c_{(i-1)m + m}): (c_1, \dots, c_{(i-1)m + j}, \dots, c_{n}) \in C_1, \]
	i.e., $C_1^{(i)}$ is obtained by deleting all the coordinates of $C_1$ except for $\{(i-1)m+1, \dots, im \}$ with $1 \leqslant i \leqslant \ell$.
	
	Similar to the proof of Theorem \ref{thm-wd}, for each $\vec{c} \in C_2$, it can be uniquely expressed by $\vec{c}_1 \oplus p \vec{c}_2$, where $\vec{c}_1 = \pi(\vec{c}) \in C_1, \vec{c}_2 \in C_1$, and $\oplus$ means the addition operation over $\mathbb{Z}_{p^2}$.
	Obviously, for each nonzero $\vec{c} \in C_2$, we have
	\[ \wt(\vec{c}) \geqslant \wt(p \vec{c}) = \wt(p \vec{c}_1) = \wt(\vec{c}_1).  \]
	
	
	Next, we will show that $C_2$ do not have other nonzero Hamming weights.
	For otherwise, let $\vec{u} = \vec{u}_1 \oplus p \vec{u}_2 \in C_2$ be the nonzero vector such that $n - \ell < \wt(\vec{u}) < n$.
	Obviously, if $\wt(\vec{u}_1) \in \{ 0, n \}$ or $\wt(\vec{u}_2) \in \{ 0, n \}$, then $\wt(\vec{u}_1) = n$.
	Hence, $\wt(\vec{u}_1) = \wt(\vec{u}_2) = n - \ell$.
	Let
	\[ \vec{u}_i = (\vec{u}_i^{(1)}, \vec{u}_i^{(2)}, \dots, \vec{u}_i^{(\ell)}), i = 1,2, 1 \leqslant j \leqslant \ell, \]
	where $u_i^{(j)}$ has length $m$.
	Similar to the case $(3).(b)$ in the proof of Theorem \ref{thm-wd}, if $\vec{u}_1^{(j)} \neq \lambda \vec{u}_2^{(j)} \pmod{p}$ for all nonzero $\lambda \in \mathbb{Z}_p$ and $1 \leqslant j \leqslant \ell$, then $\wt(\vec{u}) = n$.
	Since $n - \ell < \wt(\vec{u}) < n$, then there exists $j$ such that $\vec{u}_1^{(j)} = \lambda \vec{u}_2^{(j)} \pmod{p}$.
	Considering the codeword $\vec{v} = (-\lambda p \vec{u}_1) \oplus p \vec{u}_2 \pmod{p^2}$, we have
	\[ \wt(\vec{v}) \leqslant n - m < n - \ell, \quad \text{since}~n = m \ell > \ell^2. \]
	\begin{itemize}
		\item If $\wt(\vec{v}) = 0$, then $\vec{u}_1^{(j)} = \lambda \vec{u}_2^{(j)} \pmod{p}$ for all $j$, which implies that $\wt(\vec{u}) = n - \ell$, contradiction.
		
		\item If $\wt(\vec{v}) \neq 0$, then $\vec{v}$ is a codeword in $C_2$ satisfying $p \vec{v} = \vec{0}$ and has Hamming weight $n-\ell$ or $n$, contradiction.
	\end{itemize}
	Hence, $C_2$ is also a two-weight code.
	By the first two Pless power moments \cite[\S 7.3]{HuffmanPless}, we can calculate the related weight distribution.
\end{proof}


In Tables~\ref{tab-two-weight-2},~\ref{tab-two-weight-3}, we display selected cyclic $[n, 2, w_1]_p$ codes over $\mathbb F_p$ with one nonzero weight $w_1$ or two nonzero weights $w_1$ and $w_2$ $(w_1 < w_2)$ for WSS($d$)$^*$ primes $p$ $(7 \le p \le 167)$. In the first column,
$(n, d, (\frac{\Delta_p}{p}), p)$ denotes the code length, discriminant of $H(x)$, the Legendre symbol of $\Delta_p$ mod $p$, where $\Delta_p$ is given in Section~\ref{sub-number-theory}, and prime number, respectively. Given $p$, we display one value of $d$ although there are several values as $p$ increases.
The second and third columns are the same as those in Table~\ref{tab-two-weight}.
In the fourth column, we give a check polynomial $h(x)$ over $\mathbb F_p$ of order $p-1$ or $2p+2$ and its Hensel's lift $H(x)$ over $\mathbb Z_{p^2}$.
 In the fifth column, we distinguish the codes with MDS, NMDS, or the direct sum of the 2-fold of an $[n/2, 2, w_1/2]_p$ MDS (simplex) code denoted by 2-MDS or the 4-fold of an $[n/4, 2, w_1/4]_p$ MDS (simplex) code denoted by 4-MDS.


\medskip

The computational results in Tables~\ref{tab-two-weight-2},~\ref{tab-two-weight-3} by Magma \cite{magma}, inspired us the following results.
%{\bf{Conjecture 1}}
\begin{prop}
Let $p>5$ be a prime $\equiv 3 \pmod{4}$. Then an associated cyclic code $C_1$ from Proposition~\ref{prop-order-2p+2} is an NMDS code with parameters $[2(p+1), 2, 2p]_p$ having only one-weight.
\end{prop}

\begin{proof}
Recall from Section~\ref{subsec-irr-cyc} that $\frac{\alpha^p}{\alpha} =\alpha^{p-1}$ has order $p+1$. Then by~\cite[Theorem 1]{SZS}, the code $C_1$ is a one-weight code with the nonzero weight $2(p+1) - 2= 2p$.
 Hence it is AMDS. Moreover, since it is a one-weight code, it is the 2-fold of the simplex code $\mathcal S_1$ with parameters $[p+1, 2, p]_p$. Then the dual of the code $C_1$ contains a codeword of weight 2, say a vector $(1, 0, \cdots, 0~ |~ 1, 0, \cdots, 0)$. Hence, the dual of the code $C_1$ has parameters $[2(p+1), 2p, 2]_p$, which is also AMDS. Therefore, $C_1$ is NMDS.
\end{proof}
\begin{rem}{\rm
The length $2p+2$ is the largest possible for a near MDS code of dimension $2$ over $\F_p,$ by Proposition 5.1. of \cite{DL}. Such codes are called extremal NMDS codes.
}
\end{rem}

\begin{rem}{\rm
In the fourth column of Tables~\ref{tab-two-weight-2},~\ref{tab-two-weight-3}, $h(x)$ can be distinguished by two quadratic polynomials $h_{+1}(x)$ and $h_{-1}(x)$ depending on the value of
$(\frac{\Delta_p}{p})= \pm 1$ for a fixed prime $p$. We observe that $h_{-1}(0) = h_{+1}(1) -1$.
%This is based on our Magma calculation to compute $h(x)$.

The reason is as follows.
Let $h_{-1}(x)= (x- \alpha)(x- \alpha^p)$ as in Proposition~\ref{prop-order-2p+2}, where  $\alpha \in \F_{p^2}\setminus \F_p$ satisfies $\alpha^{p+1}=-1$. Then $h_{-1}(x)= x^2 - (\alpha + \alpha^p) + \alpha \alpha^p = x^2 -(\alpha + \alpha^p)x  -1 =  x^2 -(\alpha + \alpha^p)x + (p-1)$. Thus, $h_{-1}(0)= p-1$. On the other hand, let $h_{+1}(x) = (x-1)(x-\beta)$ where $\beta \in \F_p$ has order $p-1$ by the procedure described in~\ref{sec-order-p-1}. By expanding $h_{+1}$, we have $h_{+1}(x) = x^2 - (1 + \beta)x + \beta = x^2 + (p-1 - \beta)x + \beta $. Hence $h_{+1}(1)-1 =p-1$. Therefore, we have $h_{-1}(0) = h_{+1}(1) -1$.
}
\end{rem}

    \begin{table}[htbp]
    \center
	\caption{Cyclic $[n, 2, \{w_1, w_2\}]_p$ codes for some WSS($d$)$^*$ primes $p$ ($7 \le p \le 47$)} \label{tab-two-weight-2}
\scriptsize{
	\begin{tabular}{lllll}
		\toprule
		$(n, d, (\frac{\Delta_p}{p}) ,p)$ & $\{w_1,w_2\}$  or  & $(a_1,a_2)$ or  & $h(x)$ and & MDS/NMDS or \\
		                        &  $\{w_1\}$ & $(a_1)$ & $H(x)$ & 2-MDS, 4-MDS\\
		\toprule
		$(6,85,+1,7)$ & $\{5,6\}$ & (36,12) &   $x^2 +x +5$ & MDS\\
 & & & $x^2+29x+19$  &  \\
		$(16,649,-1,7)$ & $\{14\}$ & (48) &  $x^2 +x +6$  & NMDS, 2-MDS \\
 & & & $x^2+29x+48$  &  \\
\hline
		$(10,3,+1,11)$ & $\{9,10\}$ & (100,20) &  $x^2 +4x +6$  & MDS \\
 & & & $ x^2 + 26x + 94$  &  \\
		$(24,6,-1,11)$ & $\{22\}$ & (120) &  $x^2 +2x +10$ & NMDS, 2-MDS \\
 & & & $ x^2 + 24x + 120$  &  \\
\hline

%		$(12,10,+1,13)$ & $\{11,12\}$ & (144,24) &  $x^2 +5x +7$ & MDS\\
%& & & new h(x)   &  \\

%		$(28,426,-1,13)$ & $\{24,28\}$ & (84,84) &  $x^2 +x +12$ & 4-MDS \\
%/& & & new h(x)   &  \\

		$(16,426,+1,17)$ & $\{15,16\}$ & (256,32) &  $x^2 +10x +6$ & MDS\\
& & & $x^2 + 248x + 40$   &  \\

		$(36,193,-1,17)$ & $\{32,36\}$ & (144,144) &  $x^2 +x +16$ & 4-MDS\\
& & & $x^2 + 103x + 288$   &  \\
\hline


%		$(18,5,+1,19)$ & $\{15,18\}$ & (324,36) &  $x^2 +8x +10$ & MDS\\
% & & & new h(x)   &  \\


		$(40,43081,-1,19)$ & $\{38\}$ & (360) &  $x^2 +2x +18$ & NMDS, 2-MDS \\
& & & $x^2 + 211x + 360$ & \\
\hline

		$(22,381,+1,23)$ & $\{21,22\}$ & (484,44) &  $x^2 +11x +11$ & MDS\\
& & & $x^2 + 333x + 195$ & \\

		$(48,697,-1,23)$ & $\{46\}$ & (528) &  $x^2 +x +22$ & NMDS, 2-MDS\\
& & & $x^2 + 70x + 528$ & \\

\hline
		$(28,6821,+1,29)$ & $\{27,28\}$ & (784,56) &  $x^2 +9x +19$ & MDS\\
& & & $x^2 + 415x + 425$ & \\

		$(60,46,-1,29)$ & $\{56,60\}$ & (420,420) &  $x^2 +6x +28$ & 4-MDS\\
& & & $x^2 + 64x + 840$ & \\

\hline
		$(30,915,+1,31)$ & $\{29,30\}$ & (900,60) &  $x^2 +19x +11$ & MDS\\
& & & $x^2 + 546x + 414$ & \\

		$(64,384289,-1,31)$ & $\{62\}$ & (960) &  $x^2 +3x +30$ & NMDS, 2-MDS\\
& & & $x^2 + 623x + 960$ & \\

\hline
		$(36,7619,+1,37)$ & $\{35,36\}$ & (1296,72) &  $x^2 +23x +13$ & MDS\\
& & & $ x^2 + 874x + 494 $ & \\

		$(76,29321,-1,37)$ & $\{72,76\}$ & (684,684) &  $x^2 +x +36$ & 4-MDS\\
& & & $x^2 + 519x + 1368$ & \\

\hline
		$(40,1369205,+1,41)$ & $\{39,40\}$ & (1600,80) &  $x^2 +23x +17$ & MDS\\
& & & $x^2 + 1171x + 509$ & \\

		$(84,1523449,-1,41)$ & $\{80,84\}$ & (840,840) &  $x^2 +7x +40$ & 4-MDS\\
& & & $x^2 + 1237x + 1680$ & \\

\hline
		$(42,2022,+1,43)$ & $\{41,42\}$ & (1764,84) &  $x^2 +13x +29$ & MDS\\
& & & $x^2 + 1260x + 588$ & \\

		$(88,125563,-1,43)$ & $\{86\}$ & (1848) &  $x^2 +x +42$ & NMDS, 2-MDS \\
& & & $x^2 + 1420x + 1848$ & \\

\hline
		$(46,2085,+1,47)$ & $\{45,46\}$ & (2116,92) &  $x^2 +35x +11$ & MDS\\
& & & $x^2 + 599x + 1609$ & \\

		$(96,2554369,-1,47)$ & $\{94\}$ & (2208) &  $x^2 + 3x +46$ & NMDS, 2-MDS \\
& & & $x^2 + 1601x + 2208$ & \\

		\bottomrule
	\end{tabular}
}
\end{table}



    \begin{table}[htbp]
    \center
	\caption{Cyclic $[n, 2, \{w_1, w_2\}]_p$ codes for some WSS($d$)$^*$ primes $p$ ($53 \le p \le 97$)} \label{tab-two-weight-3}
\scriptsize{
	\begin{tabular}{lllll}
		\toprule
		$(n, d, (\frac{d}{p}), p)$ & $\{w_1,w_2\}$  or  & $(a_1,a_2)$ or  & $h(x)$ and & MDS/NMDS or \\
		                        &  $\{w_1\}$ & $(a_1)$ & $H(x)$ & 2-MDS, 4-MDS\\
		\toprule

		$(52,945867,+1,53)$ & $\{51,52\}$ & (2704,104) &  $x^2 +38x +14$ & MDS\\
& & &  $x^2 + 1946x + 862$  &\\

		$(108,200593,-1,53)$ & $\{104,108\}$ & (1404, 1404) &  $x^2 + x +52$ & 4-MDS\\
& & & $x^2 + 902x + 2808$ & \\

\hline
		$(58,1849301,+1,59)$ & $\{57,58\}$ & (3364,116) &  $x^2 +6x +52$ & MDS\\
& & &  $x^2 + 1363x + 2117$ & \\
		$(120, 5895841,-1,59)$ & $\{118\}$ & (3480) &  $x^2 + 12x +58$ & NMDS, 2-MDS \\
& & & $x^2 + 2431x + 3480$ & \\

\hline
		$(60,1210683,+1,61)$ & $\{59,60\}$ & (3600,120) &  $x^2 +6x +54$ & MDS\\
& & & $x^2 + 2202x + 1518$ & \\

		$(124,782295,-1,61)$ & $\{120, 124\}$ & (1860, 1860) &  $x^2 + 2x +60$ & 4-MDS \\
& & & $x^2 + 3540x + 3720$ & \\

\hline
		$(66,257765,+1,67)$ & $\{65,66\}$ & (4356,132) &  $x^2 +53x +13$ & MDS\\
& & & $x^2 + 1527x + 2961$ & \\

		$(136,2916193,-1,67)$ & $\{134\}$ & (4488) &  $x^2 + x +66$ & NMDS, 2-MDS \\
& & & $x^2 + 3418x + 4488$ & \\

\hline

		$(70,23218,+1,71)$ & $\{69,70\}$ & (4900,140) &  $x^2 +57x +13$ & MDS\\
& & & $x^2 + 1832x + 3208$ & \\

		$(144,2429065,-1,71)$ & $\{142\}$ & (5040) &  $x^2 + 3x +70$ & NMDS, 2-MDS \\
& & & $x^2 + 1565x + 5040$ & \\
\hline

		$(72,917,+1,73)$ & $\{71,72\}$ & (5184,144) &  $x^2 +39x +33$ & MDS\\
& & & $x^2 + 769x + 4559$ & \\

		$(148,27614737,-1,73)$ & $\{144, 148\}$ & (2664, 2664) &  $x^2 + x +72$ & 4-MDS\\
& & & $x^2 + 5257x + 5328$ & \\
\hline
		$(78,556413,+1,79)$ & $\{77,78\}$ & (6084,156) &  $x^2 +18x +60$ & MDS\\
& & & $x^2 + 3731x + 2509$ & \\


		$(160,8784985,-1,79)$ & $\{158\}$ & (6240) &  $x^2 + 5x +78$  & NMDS, 2-MDS \\
& & & $x^2 + 5930x + 6240$ & \\
\hline
		$(82,8378,+1,83)$ & $\{81,82\}$ & (6724,164) &  $x^2 +68x +14$ & MDS\\
& & & $x^2 + 400x + 6488$ & \\

		$(168,201961,-1,83)$ & $\{166\}$ & (6888) &  $x^2 + x +82$ & NMDS, 2-MDS \\
& & & $x^2 + 914x + 6888$ & \\
\hline
		$(88,4121103,+1,89)$ & $\{87,88\}$ & (7744,176) &  $x^2 +57x +31$ & MDS\\
& & & $x^2 + 4062x + 3858$ & \\

		$(180,2638645,-1,89)$ & $\{176,180\}$ & (3960, 3960) &  $x^2 + 3x +88$ & 4-MDS\\
& & & $x^2 + 6500x + 7920$ & \\
\hline
		$(96,28355,+1,97)$ & $\{95,96\}$ & (9216,192) &  $x^2 +67x +29$ & MDS\\
& & & $x^2 + 4044x + 5364$ & \\

		$(196,18186729,-1,97)$ & $\{192, 196\}$ & (4704, 4704) &  $x^2 + x +96$ & 4-MDS\\
& & & $x^2 + 4269x + 9408$ & \\
\bottomrule

%		$(100,6,+1,101)$ & $\{99,100\}$ & (10000,200) &  $x^2 +60x +40$ & MDS \\
%		$(204,8,-1,101)$ & $\{200,204\}$ & (5100, 5100) &  $x^2 + 2x +100$ & 4-MDS\\

%		$(102,8,+1,103)$ & $\{101,102\}$ & (10404,204) &  $x^2 +25x +77$ & MDS \\
%		$(208,5,-1,103)$ & $\{206\}$ & (10608) &  $x^2 + x +102$ & NMDS, 2-MDS \\

%		$(106,10,+1,107)$ & $\{105,106\}$ & (11236,212) &  $x^2 +61x +45$ & MDS \\
%		$(216,8,-1,107)$ & $\{214\}$ & (11448) &  $x^2 + 2x +106$ & NMDS, 2-MDS \\

%		$(108,7,+1,109)$ & $\{107,108\}$ & (11664,216) &  $x^2 +13x +95$ & MDS \\
%		$(220,8,-1,109)$ & $\{216, 220\}$ & (5940, 5940) &  $x^2 + 2x +108$ & 4-MDS \\

%		$(112,7,+1,113)$ & $\{111,112\}$ & (12544,224) &  $x^2 +79x +33$ & MDS \\
%		$(228,29,-1,113)$ & $\{224, 228\}$ & (6384, 6384) &  $x^2 + 5x +112$ & 4-MDS \\

%		$(126,8,+1,127)$ & $\{125,126\}$ & (15876,252) &  $x^2 +30x +96$ & MDS \\
%		$(256,5,-1,127)$ & $\{254\}$ & (16128) &  $x^2 + x + 126$ & NMDS, 2-MDS \\

%		$(130,38,+1,131)$ & $\{129,130\}$ & (16900,260) &  $x^2 +116x +14$ & MDS \\
%		$(264,8,-1,131)$ & $\{262\}$ & (17160) &  $x^2 + 2x + 130$ & NMDS, 2-MDS \\

%		$(136,7,+1,137)$ & $\{135,136\}$ & (18496,272) &  $x^2 +123x +13$ & MDS \\
%		$(276,5,-1,137)$ & $\{272, 276\}$ & (9384, 9384) &  $x^2 + x + 136$ & 4-MDS \\

%		$(138,5,+1,139)$ & $\{137,138\}$ & (19044,276) &  $x^2 +10x +128$ & MDS \\
%		$(280,8,-1,139)$ & $\{278\}$ & (19320) &  $x^2 + 2x + 138$ & NMDS, 2-MDS \\

%		$(148,7,+1,149)$ & $\{147,148\}$ & (21904,296) &  $x^2 +91x +57$ & MDS \\
%		$(300,8,-1,149)$ & $\{296,300\}$ & (11100,11100) &  $x^2 + 2x + 148$ & 4-MDS \\

%		$(150,5,+1,151)$ & $\{149,150\}$ & (22500,300) &  $x^2 +94x +56$ & MDS \\
%		$(304,13,-1,151)$ & $\{302\}$ & (22800) &  $x^2 + 3x + 150$ & NMDS, 2-MDS\\

%		$(156,3,+1,157)$ & $\{155,156\}$ & (24336,312) &  $x^2 +83x +73$ & MDS \\
%		$(316,5,-1,157)$ & $\{312,316\}$ & (12324,12324) &  $x^2 + x + 156$ & 4-MDS \\

%		$(162,10,+1,163)$ & $\{161,162\}$ & (26244,324) &  $x^2 +68x +94$ & MDS \\
%		$(326,5,-1,163)$ & $\{326\}$ & (26568) &  $x^2 + x + 162$ & NMDS \\

%		$(166,6,+1,167)$ & $\{165,166\}$ & (27556,332) &  $x^2 +136x +30$ & MDS \\
%		$(336,5,-1,167)$ & $\{334\}$ & (27888) &  $x^2 + x + 166$ & NMDS \\

%		%$(172,+1,173)$ & (171,172) & (29584,344) &  $x^2 +170x +2$ \\
%		%$(348,-1,173)$ & (344,348) & (14964,14964) &  $x^2 + x + 172$ \\
	\end{tabular}
}
\end{table}


In what follows, we describe the weight distribution of a cyclic code over  $\mathbb{Z}_{p^2}$.

Let $C = C(a,b)$ be the cyclic code of length $n$ over $\mathbb{Z}_{p^2}$ of check polynomial $h(x) = \dfrac{(1-ax)(1-bx)}{ab}$, where $n \mid p(p-1)$.
If $a \equiv b \pmod{p}$, then $C(a,b)$ has size $p^3$ by the proof of \cite[Theorem 2]{SHS},
which implies that $C$ is not free, contradiction.
Thus, we assume that $a \not\equiv b \pmod{p}$.
\begin{prop}\cite[Theorem 1]{SHS}
	Assume that $a \not\equiv b \pmod{p}$ and $e = \mathrm{ord}_p (b/a) = \mathrm{ord}_{p^2} (b/a)$.
	Then the code $C(a,b)$ has nonzero Hamming weights $\{n - \frac{n}{e}, n\}$, and their respective frequencies are
	\[ A_{n-\frac{n}{e}}(C(a,b)) = e(p^2 - 1), \quad A_n(C(a,b)) = (p^2 - 1)(p^2 - e + 1). \]
\end{prop}

\begin{coro}
	Assume that $a \not\equiv b \pmod{p}$ and $e = \mathrm{ord}_p (b/a) = \mathrm{ord}_{p^2} (b/a)$.
	Let $C_2 = C(a,b)$ and $C_1$ be the cyclic code over $\mathbb{F}_p$ with the check polynomial $\bar{h} \equiv h \pmod{p}$.
	Then $C_1$ has nonzero Hamming weights $\{n - \frac{n}{e}, n\}$, and their respective frequencies are
	\[ A_{n-\frac{n}{e}}(C_1) = e(p - 1), \quad A_n(C_1) = (p - 1)(p - e + 1). \]
	Moreover, $A_{w_i}(C_1)$ divides $A_{w_i}(C_2)$ for $i=1,2$, if and only if
	\[ (p-e+1) \mid (p+1)(p^2 - e + 1). \]
\end{coro}
\begin{proof}
	Since $C_2$ is two-Hamming-weight, then $C_1$ is two-Hamming weight and has nonzero weights $\{n - \frac{n}{e}, n\}$.
	The frequencies	of the weights can then be computed by the first two Pless power moments \cite[\S 7.3]{HuffmanPless}.
	They are the solutions of the system
	\[
	A_{n-\frac{n}{e}}(C_1) + A_n(C_1) = p^2 - 1, \quad
	\left( n - \frac{n}{e} \right) A_{n-\frac{n}{e}}(C_1) + n A_n(C_1) = pn(p-1).
	\]
	Obviously, $A_{n-\frac{n}{e}}(C_1) \mid A_{n-\frac{n}{e}}(C_2)$ always holds, while $A_{n-\frac{n}{e}}(C_1) \mid A_{n-\frac{n}{e}}(C_2)$ is true exactly when $(p-e+1) \mid (p+1)(p^2 - e + 1)$.
\end{proof}

\begin{rem}{\rm
	The code $C_2$ is projective if and only if $\mathrm{ord}_{p^2} (b/a) \geqslant n$.
	If $e = n$, then $C_2$ is MRD in the sense of \cite[Chap. 12]{SAS}, and the related $C_1$ is MDS.
}
\end{rem}


\section{Conclusion and open problem}
In this paper we have considered a family of integral second-order linear recurrences parametrized by a an integer $d>0.$
When this recurrence admits the same period modulo $p$ and $p^2$ for some prime $p,$ the quadratic field $\Q(\sqrt{d})$
is not $p$-rational, a remarkable property from an algebraic number theory standpoint. We have, for these special primes
considered the two cyclic codes over $\F_p$ and $\Z_{p^2}$
and compared their respective weight distributions. It would be interesting to use these codes to obtain relations and necessary existence conditions between $p$ and $d.$ This could have a significant impact on algebraic number theory.\\

\begin{thebibliography}{99}

%\bibitem{A} V. Andrejic, On Fibonacci powers, Univ. Beograd. Publ. Elektrotehn. Fak.Ser. Mat. 17 (2006), 38--44.

\bibitem{magma} W. Bosma, J. Cannon, C. Playoust, The Magma algebra system I: The user language, Journal of Symbolic Computation, 24(3-4):235--265, 1997.

\bibitem{B} Z. Bouazzaoui, On periods of Fibonacci sequences and real quadratic $p$-rational fields, Fibonacci Quarterly, 58(5):103--110, 2020.

\bibitem{B2} Z. Bouazzaoui, Fibonacci numbers and real quadratic $p$-rational fields, Periodica Mathematica Hungarica, 81:123--133, 2020.

\bibitem{BM} A. E. Brouwer, H. Van Maldeghem, Fragments of a text on strongly regular graphs, available at: \url{https://www.win.tue.nl/\~aeb/}

\bibitem{Co} J. Coates, $p$-adic $L$-functions and Iwasawa's theory, in: Algebraic Number Fields: $L$-Functions and Galois Properties, Proceedings of a Symposium at the University of Durham, 1975, Academic Press, London, 1977, pp. 269--353.

\bibitem{DL} S. Dodunekov, I. Landgev, On near MDS codes, Journal of Geometry, 54:30--42, 1995.

\bibitem{EJ} A-E. Elsenhans, J. Jahnel, The Fibonacci sequence modulo $p^2,$ an investigation by computer for $p < 10^{14}$, arXiv:1006.0824, 2010.

\bibitem{EGS-MDS-weight-number} M. F. Ezerman, M. Grassl, P. Sol\'e, The weights in MDS codes, IEEE Transactions on Information Theory, 57(1):392--396, 2011.

\bibitem{HuffmanPless} W. C. Huffman, V. Pless, Fundamentals of Error-Correcting Codes, Cambridge University Press, Cambridge, 2003.

\bibitem{Greenberg} R. Greenberg, Galois representations with open image, Annales math\'ematiques du Qu\'ebec, 40(1):83--119, 2016.

\bibitem{Kl} J. Klaska, Short remark on Fibonacci-Wieferich primes, Acta Mathematica Universitatis Ostraviensis, 15(1):21--25, 2007.

\bibitem{K} S. Kurz, Non-projective two-weight codes, Entropy, 26(4):289, 2024.
%\bibitem{L} E. Lehmer, On congruences involving Bernoulli numbers and quotients of Fermat and Wilson, Ann.Math. 39 (1938) 350--360.

\bibitem{LN} R. Lidl, H. Niederreiter, Finite Fields, Encyclopedia of Mathematics and its Applications, Cambridge University Press, Cambridge, 1996.

\bibitem{MacSlo} F. J. MacWilliams, N. J. A. Sloane, The Theory of Error-Correcting Codes, Elsevier-North-Holland, Amsterdam, 1977.

\bibitem{McG} G. McGuire, An approach to Hensel lemma, Irish Mathematical Society Bulletin, 47:15--21, 2001.

\bibitem{MR} R. J. McIntosh, E. L. Roettger, A search for Fibonacci-Wieferich and Wolstenholme primes, Mathematics of Computation, 76(260):2087--2094, 2007.

\bibitem{Movahhedi-Nguyen} {A. , Movahhedi}, { Q. D. Nguyen }, {Sur l'arithm\'{e}tique des corps de nombres $p$-rationnels.}
S\'{e}minaire de th\'eorie des nombres de Paris, 1987-1988, (1990), Bikh$\ddot{a}$user, Vol 81, 155-200.

\bibitem{NSW} J. Neukirch, A. Schmidt, K. Wingberg, Cohomology of Number Fields, vol. 323, Springer, 2013.

\bibitem{QS} P. Qi, M. Stokes, On the non p-rationality and Iwasawa invariants of certain real quadratic fields, arXiv:2408.03836, 2024.

\bibitem{R} M. Renault, The period, rank, and order of the $(a, b)$-Fibonacci sequence mod $m$, Mathematics Magazine, 86:372--380, 2013.

\bibitem{SAS} M. Shi, A. Alahmadi, P. Sol? Codes and Rings: Theory and Practice, Academic Press, London, 2017.

\bibitem{Se} J-P. Serre, Cohomologie galoisienne, vol. 5, Springer, 1994.

\bibitem{SHS} M. Shi, T. Helleseth, P. Sol\'e, Strongly regular graphs from reducible cyclic codes, Journal of Algebraic Combinatorics, 55:173--184, 2022.

\bibitem{SLS} M. Shi, R. Liu, P. Sol\'e, Some two-weight codes over chain rings and their strongly regular graphs, Graphs and Combinatorics, 40:84, 2024.

\bibitem{SZS} M. Shi, Z. Zhang, P. Sol\'e, Two-weight codes and second order recurrences, Chinese Journal of Electronics, 28(6):1127--1130, 2019.

\bibitem{V1} G. Vega, Determining the full weight distribution of any irreducible cyclic code over any finite field of dimension two, arXiv:1606.08510, 2016.

\bibitem{V2} G. Vega, A critical review and some remarks about one- and two-weight irreducible cyclic codes, Finite Fields and Their Applications, 33(2):1--13, 2015.

\bibitem{WS-DM} X. Wang, M. Shi, Gray images of cyclic codes over $\mathbb{Z}_{p^2}$ and $\mathbb{Z}_p \mathbb{Z}_{p^2}$, Discrete Mathematics, 346(7):113382, 2023.

\bibitem{SS}
Z. H. Sun, Z. W. Sun, Fibonacci numbers and Fermat's last theorem, Acta Arithmetica, 60(4):371--388, 1992.

\bibitem{W}
D. D. Wall, Fibonacci series modulo $m$, American Mathematical Monthly, 67(6):525--532, 1960.
\end{thebibliography}



\end{document}

We will need a Lemma
{\lem\label{binet} For any odd prime $p$ we have $F_p\equiv 5^{(p-1)/2}\pmod{p}.$}

\begin{proof}
By Binet's formula
$$F_p=\frac{(1+\sqrt{5})^p-(1-\sqrt{5})^p}{2^p \sqrt{5}}.$$
Expanding by the binomial formula we get
$$F_p=\frac{\sum_{2t+1 \le p}{p \choose 2t+1}5^t}{2^{p-1}}$$
By Fermat $2^{p-1}\equiv 1 \pmod{p}.$ It is well-known that $p|{p \choose 2t+1}$ for $2t+1<p.$
The result follows.
\end{proof}

%We are now ready for the main result.
The following is a special case of \cite[Thm 1, (iv)]{MR} with $P=1,\,Q=-1.$
{\prop If $p$ is WSS then $F_p \equiv (\frac{5}{p}) \pmod{p^2}.$ }

\begin{proof} Uses criterion 1 or 6.
\begin{itemize}
\item{$p\equiv \pm 1 \pmod{5}.$} By point 1 we know that $F_{p-1}\equiv 0 \pmod{p^2}.$ Then by point 4 we get $F_{p}\equiv 1 = (\frac{5}{p})\pmod{p^2}.$

\item{$p\equiv \pm 2 \pmod{5}.$} By point 1 we know that $F_{p+1}\equiv 0 \pmod{p^2}.$ Then by point 4 we get $F_{p-1}\equiv 1 \pmod{p^2}.$ But $F_p=F_{p+1}-F_{p-1}$ by Fibonacci recursion, implying $F_{p}\equiv -1 = (\frac{5}{p})\pmod{p^2}.$

\end{itemize}
\end{proof}

Searching for a contradiction we evaluate $F_{p}\pmod{p^2},$ by a similar argument to Lemma \ref{binet}.

{\prop If $p$ is WSS then $$F_{p}\equiv (1/2^{p-1}) ( 5^{(p-1)/2}-\frac{p}{(p-1)}-\frac{pH((p-3)/2)}{2} )\pmod{p^2}.$$ }
\begin{proof}
Note first that it is elementary by properties of primes that $$\binom{p}{k}\equiv (-1)^{k-1}\frac{p}{k} \pmod{p^2}.$$
Expanding the Binet formula yields
$$2^{p-1} F_p \equiv 5^{(p-1)/2}+\sum_{2t+1 \le p}{p \choose 2t+1}5^t,$$
or after the preceding remark

$$2^{p-1} F_p \equiv 5^{(p-1)/2}+\sum_{2t+1 \le p}\frac{p}{2t+1}.$$
The sum on the RHS is $H(p-2)-\frac{H((p-3)/2)}{2}.$
The result follows then upon using Wolstenholme theorem
to write $H(p-2)\equiv \frac{-1}{p-1} \pmod{p^2}.$
\end{proof}

So in comparison with the preceding strategy we can focus on harmonic numbers. %Can we evaluate $H((p-3)/2))\pmod{p^2}?$ Also needed is an evaluation of $2^{p-1} \pmod{p^2}.$
An old result of Lehmer \cite{L} is that
$$H(\frac{p-1}{2})\equiv -2q_p(2)+pq_p(2) \pmod{p^2},$$

where $q_p(2)=\frac{2^{p-1}-1}{p},$ a Fermat quotient.

{\coro If $p$ is WSS then $$2^pF_p\equiv 2\times 5^{(p-1)/2}-pH(\frac{p-1}{2})-4p \pmod{p^2}. $$}
\begin{proof}
By definition of the Harmonic numbers  $H(\frac{p-3}{2})=H(\frac{p-1}{2})+\frac{2}{1-p}.$
But $\frac{2}{1-p}\equiv 2+2p \pmod{p^2}.$ Thus $H(\frac{p-3}{2})\equiv H(\frac{p-1}{2})+2+2p \pmod{p^2}.$
Note also that $\frac{p}{p-1} \equiv p(p+1) \equiv p \pmod{p^2}.$ The result follows by reporting these congruences in the previous Proposition.
\end{proof}

To express the necessary condition as a function of $q_p(2)$ and $q_p(5)$ only we need the following Lemma.

{\lem\label{trick} Let $\epsilon=(\frac{5}{p})$ a Legendre symbol. Then $$ 5^{\frac{p-1}{2}}\equiv \epsilon (1+\frac{p q_p(5)}{2})\pmod{p^2}$$}

\begin{proof} By Gauss criterion $ 5^{\frac{p-1}{2}}\equiv \epsilon +ph$ for some $h$ in $\Z_{p^2}.$ Squaring this relation determines $h \pmod{p}$ after using the definition of the Fermat quotient.
\end{proof}
{\thm Assume $p>5$ is a WSS prime.

\begin{itemize}
\item if $p\equiv \pm 1 \pmod{5}$ then $q_p(5)\equiv 2 \pmod{p}$ or $$5^{p-1}\equiv 1+2p\pmod{p^2}.$$
\item if $p\equiv \pm 2 \pmod{5}$ then $q_p(5)\equiv 2q_p(2)-2 \pmod{p}$ or $$5^{p-1}\equiv 2^p-1-2p\pmod{p^2}.$$
\end{itemize}}
\begin{proof}
Let $\epsilon=(\frac{5}{p})$ a Legendre symbol. We know that $F_p\equiv \epsilon \pmod{p^2}.$ By definition of the Fermat quotient $2^p=2(1+pq_p(2)).$
The LHS of Corollary 1 is then $2\epsilon (1+pq_p(2)).$ Combining Lemma \ref{trick} with Lehmer's congruence allows us to compute the RHS. After simplifications
we get $$\epsilon (q_p(2)-q_p(5)) \equiv q_p(2)-2 \pmod{p}.$$ The result follows by discussing on the value of $\epsilon.$
\end{proof}
It is easy to construct an irreducible
 quadratic polynomial over $\F_p$ that divides $x^{(p+1)/2}-1$ for $p>3.$ Indeed the only roots common to that polynomial and to $x^{p-1}-1$ are $\pm 1.$
 This can be seen from $GCD(x^{(p+1)/2}-1,x^{p-1}-1)=1$ which results in turn from $GCD((p+1),p-1)=1.$ Therefore the only roots of
 $x^{(p+1)/2}-1$ in $\F_p$ are $\pm 1.$ If $p>3$ then $(p+1)/2>2$ and there are roots of $x^{(p+1)/2}-1$ not in $\F_p.$ Let $\alpha$ be such a root.
 Then we can take $h(x)=(x-\alpha)(x-\alpha^p).$
 By Frobenius action it is in $\F_p[x]$ and irreducible since its roots are not $\F_p.$

